\documentclass[10pt,a4paper]{amsart}
\usepackage[margin=19mm]{geometry}
\usepackage{amsmath,amssymb,mathtools}
\usepackage[scr=rsfs,cal=euler]{mathalfa}
\usepackage{xfrac}
\usepackage[unicode,hidelinks,bookmarks=false]{hyperref}
\newcommand{\ie}{i.e.\ }
\newcommand{\cf}{cf.\ }
\newcommand{\resp}{respectively\ }
\newcommand{\Subsection}{\subsection}
\providecommand{\nobreakdash}{\nobreak\hbox{-}}
\setlength{\emergencystretch}{2em}
\multlinegap0pt
\usepackage[shortlabels]{enumitem}
\usepackage{amssymb}
\usepackage{float}
\usepackage{mathtools}
\usepackage{xcolor}
\usepackage{dsfont}
\usepackage{tikz}
\newtheorem{lemma}{Lemma}
\newcommand{\NDPF}{\mathrm{PF}^{\uparrow}}
\newcommand{\x}{\mathbf{x}}
\newcommand{\y}{\mathbf{y}}
\title{The Defective Parking Space and Defective Kreweras Numbers}
\author{Rebecca E. Garcia, Pamela E. Harris, Alex Moon, Aaron Ortiz, Lauren J. Quesada, Cynthia Marie Rivera S\'anchez, Dwight Anderson Williams II, Alexander N. Wilson}
\date{Source: arXiv:2405.14635v4 (CC BY 4.0)}
\begin{document}
\maketitle

\noindent\emph{Source and licence.} The Defective Parking Space and Defective Kreweras Numbers. Version arXiv:2405.14635v4 (CC BY 4.0), distributed by arXiv under the Creative Commons Attribution 4.0 International licence, http://creativecommons.org/licenses/by/4.0/, as recorded in the arXiv licence metadata for this exact version and shown on the abstract page https://arxiv.org/abs/2405.14635v4. Full licence notice: \texttt{licenses/arxiv-2405.14635-CC-BY-4.0.txt}.\par\smallskip

\noindent\textit{Editorial scope.} Counted unit: the article's lemma lem:ndpf-bij (source0.tex lines 749--762). The article's complete original proof of it is delivered with it (source0.tex lines 764--818, one \texttt{proof} environment of 55 lines). No mathematics is written, repaired or re-derived: the statement and the proof are the article's own lines, in the article's own notation, and no retained line is substituted or re-flowed.  No further source-local context is needed: every cross-reference of every retained line resolves inside the two delivered spans.  Nothing was omitted from any retained span, so the package carries no omission record.  No retained line contains a citation, so the wrapper carries no bibliography and no bibliography entry.  No retained line contains a figure environment, an image-inclusion command or drawing code, so no image is delivered and the package declares no asset dependency.  Editorial formatting only, in addition: an AMS \texttt{amsart} preamble that restates the article's own declarations and macros used by the retained lines, namely the environments \texttt{lemma} on separate counters, together with the standard packages those lines need.  The retained environments print in this document as Lemma 1 in order of appearance, whereas the article prints the same items with the numbering of its own full document.  The complete native source of the article as distributed in the arXiv e-print for this exact version is bundled unchanged as \texttt{sources/159616/source0.tex}.
\begin{lemma}\label{lem:ndpf-bij}
    Let $k \in [n]$ and $\x=(x_1,x_2,\ldots,x_{n+1})\in \NDPF_{n+1,n+1}(x_n \leq k)$.
    Define the  map \[\psi: \NDPF_{n+1,n+1}(x_{n+1}\leq k) \to M_{n}(x_n \leq k)\]
    by 
    \[\psi(\x )=(\x',i),\] 
    where 
    $i = \max(F(\x))$ and
    \[\x' = (x_1,x_2,\ldots,x_{i-1},x_{i+1},x_{i+2},\ldots,x_{n+1}).\]
Then 
$\psi$ is a bijection.
Moreover, 
    $F(\x') = \{j \in F(\x) \mid j \leq i\}$.

\end{lemma}

\begin{proof} Fix $\x=(x_1,x_2,\ldots,x_{n+1}) \in \NDPF_{n+1,n+1}(x_{n+1} \leq k)$ and  let $i=\max(F(\x))$. Let \[\x'=(x'_1,x'_2,\dots,x'_n)=(x_1,x_2,\ldots,x_{i-1},x_{i+1},x_{i+2},\ldots,x_{n+1}),\] be the first component of $\psi(\x)$.


    \begin{itemize}[leftmargin=.2in]
        \item Well-defined: 
        Since $\x$ is a nondecreasing parking function, we know that $x_i\leq i$ for all $i\in[n+1]$.
        Moreover, since $i=\max(F(\x))$, we know that $x_\ell<\ell$ for all $i<\ell\leq n+1$. 
        Now removing $x_i$ from $\x$ to construct $\x'$ ensures that $\x'$
        satisfies the following: 
        \begin{itemize}
            \item If $1\leq j < i$, then $x'_j = x_j\leq j$, and 
            \item if $i\leq j\leq n$, then $x'_j = x_{j+1}< j+1$, so $x_j'\leq j$.
        \end{itemize}
        The above two condition ensure that $\x'\in\NDPF_{n,n}$. 
        By assumption $x_{n+1}\leq k$, so $x'_n=x_{n+1}\leq k$, which implies that $\x'\in\NDPF_{n,n}(x_n'\leq k)$, as desired. 
        Then $\psi$ is well-defined on the first component. 
        
        Because $k\leq n$, we have that $x_{n+1} < n+1$ is not a fixed point of $\x$, so $i < n+1$. As $i=\max(F(\x))$, we know $x_{i+1}$ is not a fixed point of $\x$, so $x_{i+1} < i+1$. Hence, \begin{equation}\label{eq:fp-stays}i = x_i \leq x_{i+1} < i+1,\end{equation}
        so $x'_i = x_{i+1} = i$. This implies that $i$ is in $F(\x')$, so $\psi$ is well-defined on the second component.
    
        \item Surjective: Suppose $(\x', i) \in M_n$, where $\x'=(x_1',x_2',\ldots,x_n')\in\NDPF_{n,n}$ with $x_n'\leq k$. Then let
        \[\x =(x_1,x_2,\ldots,x_{n+1})= (x'_1, x'_2, \ldots, x'_{i-1}, i, x'_{i}, x'_{i+1}, \ldots, x'_{n}).\]
        We show that $\x\in\NDPF_{n+1,n+1}(x_{n+1}\leq k)$ and $\psi(\x) = (\x', i)$. We proceed by showing that $x_j\leq j$ for all $j\in[n+1]$.
        In the case that $1\leq j < i$ we have $x_j = x'_j \leq j$. 
        When $j = i$ we have $x_i = i$.
        When $i<j\leq n+1$ we have $x_j = x'_{j-1} \leq j-1 \leq j$.
        Thus $x_j\leq j$ for all $j\in[n+1]$, so $\x$ is a parking function. Because $x'_i = i$, we have
        \[x_{i-1} < x_i = x_{i+1},\]
        so $\x$ is nondecreasing. Then $\x \in \NDPF_{n+1,n+1}$. Because $x_j = x'_{j-1} \leq j-1 < j$, for $j>i$, $i$ must be the index of the last fixed point of $\x$, so $\psi(\x) = (\x', i)$ as desired.

        \item Injective: 
         Suppose $\x=(x_1,x_2,\ldots,x_{n+1}), \y=(y_1,y_2,\ldots,y_{n+1}) \in \NDPF_{n+1,n+1}$ and assume that $\psi(\x)=\psi(\y)$. 
         This implies that  $\max(F(\x))=\max(F(\y))$, which we denote by $i$. 
         Moreover, the first coordinates are equal, i.e. $\x'=\y'$, which implies that $x'_j = y'_j$ for all $j\in[n+1]\setminus\{i\}$. By the definition of $\psi$, this implies that $x_j = y_j$ for all $j\in[n+1]\setminus\{i\}$. Then as $i=\max(F(\x))=\max(F(\y))$, we have that $x_i=i=y_i$. Thus $\x=\y$, and $\psi$ is injective.
    \end{itemize}
    Hence, $\psi$ is a bijection. 
    
    To conclude, we prove that $F(\x) = \{j \in F(\x') \mid j \leq i\}$ by establishing two set containments:
    \begin{itemize}[leftmargin=.2in]
        \item $F(\x) \subseteq \{j \in F(\x') \mid j \leq i\}$: Suppose $j \in F(\x)$. Recall $i=\max(F(\x))$. Hence,
        \[x'_j = \begin{cases}
            x_j & \mbox{if }1\leq j < i\\
            x_{j+1} & \mbox{if }j=i.
        \end{cases}\]
        If $1\leq j < i$, then  $x'_j =x_j < j$. 
        If $j = i$, then from equation \eqref{eq:fp-stays} we know that $x_{j+1} = j$, so $j \in F(\x')$. Thus $j\in\{F(\x')\mid j\leq i\}$.
        \item $ \{j \in F(\x') \mid j \leq i\}\subseteq F(\x)$: Suppose $j \in F(\x')$ and $1\leq j \leq i$. Then we have
        \[x_j = \begin{cases}
            x'_j &\mbox{if } 1\leq j < i\\
            i & \mbox{if }j=i.
        \end{cases}\]
        Hence, if $1\leq j< i$, then $\delta(\x)_j=x_j-j=x'_j-j=0$ and $j\in F(\x)$. If $j=i$, then $\delta(\x)_j=x_j-j=i-i=0$ and so $j\in F(\x)$. 
        In both cases we get that         $j \in F(\x)$, as desired.\qedhere
    \end{itemize}
\end{proof}

\end{document}
